% GENERATED FILE. Edit canonical lessons, then run npm run generate:books.
% canonical-source-hash: fnv1a64:7822e77b7513f8a3
% canonical-lessons: LA-C554-subsellium, LA-C554-urceus, LA-C554-repono, LA-C554-sarcio, LA-C554-pecto, LA-R554-review-553-554, LA-R554-review-final-554

\chapter{Household chores}
\label{ch:la-household-chores}
\hlchaptermodality{554}

\begin{chapteropening}
\textbf{By the end of this chapter:} \emph{I can talk about chores at home.}

Complete the last lesson of chapter 554: I can talk about chores at home.
\end{chapteropening}

\section[subsellium, subselliī]{subsellium, subselliī — a bench, low seat}

\label{lesson:LA-C554-subsellium}



\begin{quote}
Before the new one: say the Latin for a hinge, then the Latin for provisions.
\end{quote}

\subsection*{You'll want to know: subsellium, subselliī}
\textbf{subsellium, subselliī} — \textquotedblleft{}a bench, low seat\textquotedblright{}.

\begin{quote}
Discipulī in subselliīs sedent.
\end{quote}

\begin{quote}
\emph{The pupils sit on benches.}
\end{quote}

Neuter, second declension: \emph{sub} + \emph{sella} — lower than a chair. Also the benches of a law court.

\begin{cousinweb}
\textbf{subsellium} is built on \emph{sub} \textquotedblleft{}below\textquotedblright{} and \emph{sella} \textquotedblleft{}seat, chair\textquotedblright{} (Lewis \& Short).
\end{cousinweb}

\begin{grammarlens}[title={What you've built}]
One more word for telling, explaining and describing at B1.
\end{grammarlens}

\subsection*{Your turn}
\begin{itemize}
\raggedright
  \item \emph{Say it:} \emph{subsellium, subselliī}
  \item \emph{Say it:} \emph{subsellium, subselliī}, once more
  \item \emph{Say it:} a sentence of your own with \emph{subsellium, subselliī}
  \item \emph{Recall:} say the Latin for a hinge, then the Latin for provisions, then say \emph{subsellium, subselliī} again
\end{itemize}

\begin{tcolorbox}[breakable,colback=teal!4,colframe=teal!35!black,title={Before you move on}]
What does subsellium, subselliī mean? (\textquotedblleft{}A bench, low seat\textquotedblright{}.) Say it once more.
\end{tcolorbox}
\section[urceus, urceī]{urceus, urceī — a jug, pitcher}

\label{lesson:LA-C554-urceus}



\begin{quote}
Before the new one: say the Latin for provisions, then the Latin for a bench.
\end{quote}

\subsection*{You'll want to know: urceus, urceī}
\textbf{urceus, urceī} — \textquotedblleft{}a jug, pitcher\textquotedblright{}.

\begin{quote}
Ancilla aquam urceō portat.
\end{quote}

\begin{quote}
\emph{The maid carries water in a jug.}
\end{quote}

Masculine, second declension.

\begin{cousinweb}
\textbf{urceus} is a word of uncertain origin; Lewis \& Short compares Greek \emph{hurche} \textquotedblleft{}a jar\textquotedblright{}, and EtymDB links it with \emph{orca} \textquotedblleft{}a cask\textquotedblright{}, but nothing is settled.
\end{cousinweb}

\begin{grammarlens}[title={What you've built}]
One more word for telling, explaining and describing at B1.
\end{grammarlens}

\subsection*{Your turn}
\begin{itemize}
\raggedright
  \item \emph{Say it:} \emph{urceus, urceī}
  \item \emph{Say it:} \emph{urceus, urceī}, once more
  \item \emph{Say it:} a sentence of your own with \emph{urceus, urceī}
  \item \emph{Recall:} say the Latin for provisions, then the Latin for a bench, then say \emph{urceus, urceī} again
\end{itemize}

\begin{tcolorbox}[breakable,colback=teal!4,colframe=teal!35!black,title={Before you move on}]
What does urceus, urceī mean? (\textquotedblleft{}A jug, pitcher\textquotedblright{}.) Say it once more.
\end{tcolorbox}
\section[repōnō, repōnere]{repōnō, repōnere — to put back; to store away}

\label{lesson:LA-C554-repono}



\begin{quote}
Before the new one: say the Latin for a bench, then the Latin for a jug.
\end{quote}

\subsection*{You'll want to know: repōnō, repōnere}
\textbf{repōnō, repōnere} — \textquotedblleft{}to put back; to store away\textquotedblright{}.

\begin{quote}
Post cēnam ancilla vāsa in armāriō repōnit.
\end{quote}

\begin{quote}
\emph{After dinner the maid puts the dishes back in the cupboard.}
\end{quote}

\emph{repōnō, repōnere, reposuī, repositum}. English \emph{repository}.

\begin{cousinweb}
\textbf{repōnō} is built on \emph{re-} \textquotedblleft{}back\textquotedblright{} and \emph{pōnō} \textquotedblleft{}to put\textquotedblright{} (EtymDB; Lewis \& Short). English \emph{repository} comes from it.
\end{cousinweb}

\begin{grammarlens}[title={What you've built}]
One more word for telling, explaining and describing at B1.
\end{grammarlens}

\subsection*{Your turn}
\begin{itemize}
\raggedright
  \item \emph{Say it:} \emph{repōnō, repōnere}
  \item \emph{Say it:} \emph{repōnō, repōnere}, once more
  \item \emph{Say it:} a sentence of your own with \emph{repōnō, repōnere}
  \item \emph{Recall:} say the Latin for a bench, then the Latin for a jug, then say \emph{repōnō, repōnere} again
\end{itemize}

\begin{tcolorbox}[breakable,colback=teal!4,colframe=teal!35!black,title={Before you move on}]
What does repōnō, repōnere mean? (\textquotedblleft{}To put back; to store away\textquotedblright{}.) Say it once more.
\end{tcolorbox}
\section[sarciō, sarcīre]{sarciō, sarcīre — to mend, patch}

\label{lesson:LA-C554-sarcio}



\begin{quote}
Before the new one: say the Latin for a jug, then the Latin for to put back.
\end{quote}

\subsection*{You'll want to know: sarciō, sarcīre}
\textbf{sarciō, sarcīre} — \textquotedblleft{}to mend, patch\textquotedblright{}.

\begin{quote}
Avia tunicam nepōtis sarcit.
\end{quote}

\begin{quote}
\emph{Grandmother mends her grandson's tunic.}
\end{quote}

\emph{sarciō, sarcīre, sarsī, sartum}. Of clothes, roofs and friendships.

\begin{cousinweb}
\textbf{sarciō} is inherited from a Proto-Indo-European root meaning \textquotedblleft{}to mend, make whole\textquotedblright{} (EtymDB), though Lewis \& Short calls it doubtful.
\end{cousinweb}

\begin{grammarlens}[title={What you've built}]
One more word for telling, explaining and describing at B1.
\end{grammarlens}

\subsection*{Your turn}
\begin{itemize}
\raggedright
  \item \emph{Say it:} \emph{sarciō, sarcīre}
  \item \emph{Say it:} \emph{sarciō, sarcīre}, once more
  \item \emph{Say it:} a sentence of your own with \emph{sarciō, sarcīre}
  \item \emph{Recall:} say the Latin for a jug, then the Latin for to put back, then say \emph{sarciō, sarcīre} again
\end{itemize}

\begin{tcolorbox}[breakable,colback=teal!4,colframe=teal!35!black,title={Before you move on}]
What does sarciō, sarcīre mean? (\textquotedblleft{}To mend, patch\textquotedblright{}.) Say it once more.
\end{tcolorbox}
\section[pectō, pectere]{pectō, pectere — to comb}

\label{lesson:LA-C554-pecto}



\begin{quote}
Before the new one: say the Latin for to put back, then the Latin for to mend.
\end{quote}

\subsection*{You'll want to know: pectō, pectere}
\textbf{pectō, pectere} — \textquotedblleft{}to comb\textquotedblright{}.

\begin{quote}
Ancilla crīnēs dominae pectit.
\end{quote}

\begin{quote}
\emph{The maid combs her mistress's hair.}
\end{quote}

\emph{pectō, pectere, pexī, pexum}. \emph{Pecten}, already known, is the comb.

\begin{cousinweb}
\textbf{pectō} is inherited from a Proto-Indo-European root meaning \textquotedblleft{}to comb, shear\textquotedblright{}, and Greek \emph{peko} is a cousin (EtymDB; Lewis \& Short). \emph{pecten} \textquotedblleft{}comb\textquotedblright{} is built on it.
\end{cousinweb}

\begin{grammarlens}[title={What you've built}]
One more word for telling, explaining and describing at B1.
\end{grammarlens}

\subsection*{Your turn}
\begin{itemize}
\raggedright
  \item \emph{Say it:} \emph{pectō, pectere}
  \item \emph{Say it:} \emph{pectō, pectere}, once more
  \item \emph{Say it:} a sentence of your own with \emph{pectō, pectere}
  \item \emph{Recall:} say the Latin for to put back, then the Latin for to mend, then say \emph{pectō, pectere} again
\end{itemize}

\begin{tcolorbox}[breakable,colback=teal!4,colframe=teal!35!black,title={Before you move on}]
What does pectō, pectere mean? (\textquotedblleft{}To comb\textquotedblright{}.) Say it once more.
\end{tcolorbox}
\section[(dialogue)]{Home, books and the arts}

\label{lesson:LA-R554-review-553-554}



\begin{quote}
Say the Latin for a forecourt and a door bolt.
\end{quote}

\subsection*{Your turn}
Say these aloud:

\begin{itemize}
\raggedright
  \item a forecourt, a door bolt, a threshold, a hinge, provisions
  \item a bench, a jug, to put back, to mend, to comb
\end{itemize}

\begin{tcolorbox}[breakable,colback=teal!4,colframe=teal!35!black,title={Before you move on}]
A forecourt? (\textbf{vestibulum, vestibulī}.) A door bolt? (\textbf{pessulus, pessulī}.) A threshold? (\textbf{līmen, līminis}.) A hinge? (\textbf{cardō, cardinis}.) Provisions? (\textbf{penus, penūs}.) A bench? (\textbf{subsellium, subselliī}.) A jug? (\textbf{urceus, urceī}.) To put back? (\textbf{repōnō, repōnere}.) To mend? (\textbf{sarciō, sarcīre}.) To comb? (\textbf{pectō, pectere}.)
\end{tcolorbox}
\section[(dialogue)]{Household chores — once more}

\label{lesson:LA-R554-review-final-554}



\begin{quote}
Say the Latin for a bench and a jug.
\end{quote}

\subsection*{Your turn}
\begin{itemize}
\raggedright
  \item \emph{Say it:} a bench, a jug, to put back, to mend, to comb
\end{itemize}

\begin{tcolorbox}[breakable,colback=teal!4,colframe=teal!35!black,title={Before you move on}]
A bench? (\textbf{subsellium, subselliī}.) A jug? (\textbf{urceus, urceī}.) To put back? (\textbf{repōnō, repōnere}.) To mend? (\textbf{sarciō, sarcīre}.) To comb? (\textbf{pectō, pectere}.)
\end{tcolorbox}
